% !TeX root = ../main.tex
\section{Remaining integration and qualification priorities}
\label{sec:tracker-review}
The integrated geometry, constituent allocation, identifiers, conversion and
bounded transport checks provide the common simulation baseline. Engineering
qualification and realistic reconstruction remain open. \Cref{tab:tracker-review}
summarizes the work needed before an accepted detector release.

\begin{table}[!htbp]
\centering\small
\caption{Remaining qualification priorities for the current tracker.}
\label{tab:tracker-review}
\begin{tabular}{@{}>{\raggedright\arraybackslash}p{34mm}>{\raggedright\arraybackslash}p{106mm}@{}}\toprule
Interface & Required review outcome\\\midrule
Acceptance and response & Extend origin-ray and truth-fit studies to the luminous
region, both charges and realistic response; establish pattern-recognition
performance on the current geometry.\\
Common mechanics & Qualify global cages, joints, petal keys and thermal motion;
confirm long-stave elastic properties and handling/installation loads.\\
Cooling & Measure cold contacts and realistic power; qualify the complete
short-strip thermal path, its tight adverse margin and two-phase hydraulics,
and the long-petal lateral heat path.\\
Power and routing & Qualify ASICs, the short-strip 8 mm harness and fault domain,
long-strip adverse packing/return drop, and pixel trunk capacity, with connectors,
strain relief and installation access.\\
Material and alignment & Resolve the Geant4/TGeo graphite-interface discrepancy;
validate material locality and curved transport beyond azimuth-averaged maps.
Qualify alignment covariances and stereo-pair measurement transport.\\\bottomrule
\end{tabular}
\end{table}

Model assumptions and unqualified hardware remain explicit. Numerical checks,
publication and technical review do not grant engineering sign-off or detector
acceptance. Exact source identities and reproduction recipes accompany the
report separately from this baseline description.
